\chapter{Overlapping Generations}\label{ch:b2:overlapping-generations}

Every economy holds several generations at once: people who work, people who have retired, and children. The balance
between them is changing fast. The chart shows the number of people aged 65 and over for every 100 people of working age.

\begin{openingchart}{People aged 65 and over per 100 people aged 15--64, 1960--2025: Japan, Italy, Korea, United States,
  Nigeria and the world.}{World Bank, World Development Indicators (SP.POP.DPND.OL)}
\begin{tikzpicture}
\begin{axis}[omrecord, width=0.84\linewidth, height=0.44\linewidth, xmin=1960, xmax=2025, ymin=0, ymax=55,
    xlabel={Year}, ylabel={Per 100 of working age}, xtick={1960,1970,1980,1990,2000,2010,2020},
    xticklabel style={/pgf/number format/1000 sep=}, legend style={at={(0.5,-0.3)}, anchor=north,
    legend columns=3, /tikz/every even column/.append style={column sep=8pt}}]
  \addplot[omDef, thick] table[x=year, y=JPN] {figdata/bachelor-2/22-overlapping-generations/out/oadr.csv};
  \addplot[omThm, thick, dashed] table[x=year, y=ITA] {figdata/bachelor-2/22-overlapping-generations/out/oadr.csv};
  \addplot[omMeth, thick, dotted] table[x=year, y=KOR] {figdata/bachelor-2/22-overlapping-generations/out/oadr.csv};
  \addplot[black!70, thick, dashdotted] table[x=year, y=USA] {figdata/bachelor-2/22-overlapping-generations/out/oadr.csv};
  \addplot[omProp, thick] table[x=year, y=NGA] {figdata/bachelor-2/22-overlapping-generations/out/oadr.csv};
  \addplot[black, very thick] table[x=year, y=WLD] {figdata/bachelor-2/22-overlapping-generations/out/oadr.csv};
  \legend{Japan, Italy, Korea, United States, Nigeria, World}
\end{axis}
\end{tikzpicture}
\end{openingchart}

% figdata: oadr_summary (JPN 8.9 -> 51; ITA 14.5 -> 39.7; KOR 5.6 -> 29.3; USA 14.7 -> 28.5; NGA 6 -> 5.4; WLD 8.6 -> 16)
In 1960 Japan had 8.9 people of 65 and over for every 100 of working age; in 2025 it had 51.0. Korea went from 5.6 to 29.3, the world from 8.6 to 16.0, while
Nigeria stayed near 5. Who pays for the old, and how saving for old age shapes the capital stock, are the questions of
this chapter. The model that answers them has people who live two periods, young then old, with a new generation born
each period.

\begin{recall}
From the Year 1 volume: the Solow model, its steady state and the golden rule. From \cref{ch:b2:intertemporal-choice}: the two-period
consumption choice. From \cref{ch:b2:government-debt}: debt dynamics and the role of $i - \gamma$.
\end{recall}

\begin{notation}[Symbols of this chapter]\label{not:b2:overlapping-generations:symbols}
A period is a generation, about thirty years. $L_t$ is the number of young people in period $t$, growing at rate $n$ per period; $k_t$
capital per young person; $w_t$ the wage and $R_t$ the gross return on capital; $c_{1t}$, $c_{2t+1}$ consumption when young and old; $d$ a pension contribution.
\end{notation}

\section{The two-period model}\label{sec:b2:overlapping-generations:model}

\begin{definition}[Overlapping-generations model]\label{def:b2:overlapping-generations:olg}
An \emph{overlapping-generations model}\index{overlapping-generations model} is a model in which each person lives a
finite number of periods and a new generation is born every period, so that in each period generations of different ages
live side by side and trade.
\end{definition}

\begin{omfigure}
\begin{tikzpicture}[econ, x=2.9cm, y=0.75cm]
  \foreach \t/\lab in {0/{$t-1$}, 1/{$t$}, 2/{$t+1$}, 3/{$t+2$}} {\node[font=\footnotesize] at (\t+0.5,4.6) {period \lab};}
  \foreach \g/\y in {0/3.5, 1/2.5, 2/1.5} {
    \draw[fill=omDef!15, draw=omDef] (\g,\y-0.35) rectangle (\g+1,\y+0.35);
    \node[font=\scriptsize] at (\g+0.5,\y) {young: works, saves};
    \draw[fill=omThm!15, draw=omThm] (\g+1,\y-0.35) rectangle (\g+2,\y+0.35);
    \node[font=\scriptsize] at (\g+1.5,\y) {old: consumes};
  }
  \node[font=\footnotesize, anchor=east] at (-0.05,3.5) {born $t-1$};
  \node[font=\footnotesize, anchor=east] at (-0.05,2.5) {born $t$};
  \node[font=\footnotesize, anchor=east] at (0.95,1.5) {born $t+1$};
  \draw[guide] (1,0.8) -- (1,4.3);
  \draw[guide] (2,0.8) -- (2,4.3);
\end{tikzpicture}

{\small Generations overlap. In period $t$ the generation born at $t$ works and saves, while the generation born at $t - 1$ lives on its savings. The capital
of period $t + 1$ is what the young of period $t$ saved.}
\end{omfigure}

A person born at $t$ works when young, earns $w_t$, consumes $c_{1t}$ and saves $s_t = w_t - c_{1t}$; when old she consumes $c_{2t+1} = R_{t+1}s_t$. With $u = \ln c_1 + \beta\ln c_2$, the \omterm{def:b2:intertemporal-choice:euler-equation}{Euler equation} of \cref{ch:b2:intertemporal-choice} gives
$c_{2t+1} = \beta R_{t+1}c_{1t}$, and so $s_t = \frac{\beta}{1 + \beta}w_t$: with logarithmic utility, saving is a fixed share of the wage. Firms produce $Y = AK^\alpha L^{1-\alpha}$, so per young person $y = Ak^\alpha$, the wage is $w = (1 - \alpha)Ak^\alpha$ and the
gross return is $R = \alpha Ak^{\alpha - 1}$; capital lasts one period. Next period's capital is this period's saving, shared among $1 + n$ times as many young:
\begin{equation}\label{eq:b2:overlapping-generations:k}
k_{t+1} = \frac{s_t}{1 + n} = \frac{\beta(1 - \alpha)A}{(1 + \beta)(1 + n)}\,k_t^\alpha.
\end{equation}

\section{The steady state and its stability}\label{sec:b2:overlapping-generations:steady}

Setting $k_{t+1} = k_t = k^*$ in \eqref{eq:b2:overlapping-generations:k} gives $k^* = \left(\frac{\beta(1 - \alpha)A}{(1 + \beta)(1 + n)}\right)^{1/(1-\alpha)}$. Take $\alpha = 0.3$, $\beta = 0.6$, $n = 0.2$ and $A = 10$, an economy whose population grows 20 per cent a
generation. The coefficient is $0.6 \times 0.7 \times 10/(1.6 \times 1.2) = 2.1875$, and $k^* = 2.1875^{1/0.7} = 3.06$.

\begin{omfigure}
\begin{tikzpicture}[econ]
\begin{axis}[omrecord, grid=none, xmin=0, xmax=5, ymin=0, ymax=5, width=0.5\linewidth, height=0.5\linewidth,
    xlabel={$k_t$}, ylabel={$k_{t+1}$}, xtick={0.5,3.059}, xticklabels={$k_0$,$k^*$}, ytick=\empty]
  \addplot[black!50, domain=0:5, samples=2] {x};
  \addplot[omDef, very thick, domain=0.001:5, samples=100] {2.1875*x^0.3};
  \addplot[omThm, thick, const plot] coordinates {(0.5,0) (0.5,1.777) (1.777,1.777) (1.777,2.599) (2.599,2.599) (2.599,2.913) (2.913,2.913) (2.913,3.014)};
  \node[eqpoint] at (axis cs:3.059,3.059) {};
  \node[anchor=west, font=\footnotesize] at (axis cs:4.2,4.0) {$45^\circ$};
\end{axis}
\end{tikzpicture}

{\small Capital per young person next period against this period. The curve is concave and starts steep, so it crosses the
$45^\circ$ line once at a positive $k^*$; from any positive start, capital climbs (or falls) toward $k^*$.}
\end{omfigure}

The steady state is stable: the slope of \eqref{eq:b2:overlapping-generations:k} at $k^*$ is $\alpha < 1$, so deviations shrink by the factor $\alpha$ each generation.

\begin{model}{The overlapping-generations model}
\assumes Finitely lived households that save for their own old age, no bequests, competitive firms and markets; here two
periods of life, logarithmic utility and Cobb--Douglas production.
\predicts Saving and capital depend on the wage of the young and on the age structure; capital may exceed the golden-rule
level; \omterm{def:b2:overlapping-generations:payg}{pay-as-you-go pensions} and government debt lower the capital stock.
\tested The test of dynamic efficiency: in the United States and other large economies, capital earns more than is
invested (Abel, Mankiw, Summers and Zeckhauser 1989), and so they do not have too much capital.
\breaks Two periods of life are a coarse picture of saving over a lifetime, and households that leave bequests link the
generations, which the model leaves out.
\end{model}

\section{Too much capital is possible}\label{sec:b2:overlapping-generations:inefficiency}

In the Solow model of the Year 1 volume, the golden rule maximises steady-state consumption; with capital that lasts one
period, consumption per young person in the steady state is $Ak^\alpha - (1 + n)k$, maximised where $R = \alpha Ak^{\alpha - 1} = 1 + n$. In the example, the golden-rule stock is $(3/1.2)^{1/0.7} = 3.70$, above the steady
state, and the return $R^* = 1.371$ exceeds $1 + n = 1.2$.

\begin{definition}[Dynamic inefficiency]\label{def:b2:overlapping-generations:dynamic-inefficiency}
An economy is \emph{dynamically inefficient}\index{dynamic inefficiency} if its capital stock exceeds the golden-rule
level, so that consuming some of the capital and investing less would raise consumption in every period.
\end{definition}

In the competitive overlapping-generations economy nothing rules this out. In the steady state $R^*/(1 + n) = \alpha(1 + \beta)/(\beta(1 - \alpha))$, which is below 1 when households are
patient enough: with $\alpha = 0.3$ and $\beta = 1$ it is $0.6/0.7 = 0.857$. The young must save for their old age, and the only asset is capital, so they may accumulate too much of
it. Then a transfer from each young generation to the old one makes every generation better off.

\begin{record}[Is the capital stock too large?]
% ledger: b2:abel1989
% figdata: eff_summary (n 129, n_above 128, below CHN; median cap 47.8, inv 23.0; CHN 42.2 vs 42.7; USA 40.8 vs 21.7)
Interest rates are a poor guide, because they differ across assets and include risk. Abel, Mankiw, Summers and Zeckhauser
(1989) proposed a test that needs only the national accounts: an economy is efficient if the income capital generates
exceeds what is invested in it, year after year, so that capital is a source of resources for consumption rather than a
drain. They found the United States and the other large economies they studied to be efficient. The chart applies the
test, in one year, to 129 economies: capital income exceeds gross investment in 128 of them. China, which invested 42.7 per
cent of its output against capital income of 42.2 per cent, is the one exception.
\end{record}

\begin{omchart}{Capital income (one minus the labour share) against gross investment, per cent of GDP, both at current
  national prices, 129 economies, 2019; the $45^\circ$ line marks equality. For some countries the Penn World Table
  estimates the labour share from incomplete data.}{Penn World Table (labsh); World Bank, World Development Indicators
  (NE.GDI.TOTL.ZS); our computation}
\begin{tikzpicture}
\begin{axis}[omrecord, width=0.6\linewidth, height=0.52\linewidth, xmin=0, xmax=70, ymin=0, ymax=100,
    xlabel={Gross investment, per cent of GDP}, ylabel={Capital income, per cent of GDP}]
  \addplot[only marks, mark=*, mark size=1.3pt, omDef, fill opacity=0.6]
    table[x=inv, y=cap] {figdata/bachelor-2/22-overlapping-generations/out/efficiency.csv};
  \addplot[black!60, dashed, domain=0:70, samples=2] {x};
  \node[font=\footnotesize, anchor=west] at (axis cs:44,40) {CHN};
\end{axis}
\end{tikzpicture}
\end{omchart}

\section{Pay-as-you-go and funded pensions}\label{sec:b2:overlapping-generations:pensions}

\begin{definition}[Pay-as-you-go pension]\label{def:b2:overlapping-generations:payg}
A \emph{pay-as-you-go pension}\index{pay-as-you-go pension} system pays the pensions of the old out of the contributions
of the young in the same period.
\end{definition}

\begin{definition}[Funded pension]\label{def:b2:overlapping-generations:funded}
A \emph{funded pension}\index{funded pension} system invests each generation's contributions in assets and pays each
generation's pensions out of those assets and their returns.
\end{definition}

\begin{omfigure}
\begin{tikzpicture}[econ]
  \draw[fill=omDef!15, draw=omDef] (0,0) rectangle (3,1.2);
  \node[align=center, font=\footnotesize] at (1.5,0.6) {$L_t$ young\\each pays $d$};
  \draw[fill=black!5] (4,0) rectangle (7.2,1.2);
  \node[align=center, font=\footnotesize] at (5.6,0.6) {pension system\\$dL_t$ in, $dL_t$ out};
  \draw[fill=omThm!15, draw=omThm] (8.2,0) rectangle (12.2,1.2);
  \node[align=center, font=\footnotesize] at (10.2,0.6) {$L_{t-1} = L_t/(1 + n)$ old\\each gets $(1 + n)d$};
  \draw[->, thick] (3,0.6) -- (4,0.6);
  \draw[->, thick] (7.2,0.6) -- (8.2,0.6);
\end{tikzpicture}

{\small The budget of a pay-as-you-go system in one period. What the young pay in, the old take out the same period; since
there are $1 + n$ young people for each old one, each pensioner receives $(1 + n)d$. If the wage per worker also grows, and contributions
with it, the return is $1 + n$ times the growth factor of the wage.}
\end{omfigure}

A pay-as-you-go system pays a return of $1 + n$ on each contribution, a funded system the return on capital $R$. In a dynamically efficient economy $R > 1 + n$, so later
generations do better with a funded system. But the first old generation of a pay-as-you-go system receives pensions without
having contributed: a windfall that every later generation pays for. And because the young save less for their old age,
the capital stock falls. The listing finds the steady state of the example with and without a contribution of 1 (thousand
dalens per young person, against a wage of 9.79).

\omcode{code/bachelor-2/22-overlapping-generations/listings/olg.py}{1}{38}{The steady state of the \omterm{def:b2:overlapping-generations:olg}{overlapping-generations
model} with and without a pay-as-you-go contribution.}\label{lst:b2:overlapping-generations:olg}

\omoutput{code/bachelor-2/22-overlapping-generations/outputs/olg.txt}{A contribution of about a tenth of the wage lowers
capital per young person by a third.}\label{out:b2:overlapping-generations:olg}

% code: listings/olg.py
The contribution cuts capital per young person from 3.06 to 2.06, lowers the wage from 9.79 to 8.69 and raises the return on capital from 1.37 to 1.81. A funded system of the same size leaves
capital unchanged: what the young pay in is invested in capital, and households reduce their own saving one for one.
Moving from a pay-as-you-go system to a funded one is not free either: the generation in the middle must pay for the
pensions of its parents and save for its own, so the switch gives up the first generation's windfall only by charging
some later generation for it.

\section{The arithmetic of ageing}\label{sec:b2:overlapping-generations:ageing}

\begin{definition}[Replacement rate]\label{def:b2:overlapping-generations:replacement-rate}
The \emph{replacement rate}\index{replacement rate} of a pension system is the average pension divided by the average
wage.
\end{definition}

\begin{definition}[Old-age dependency ratio]\label{def:b2:overlapping-generations:oadr}
The \emph{old-age dependency ratio}\index{old-age dependency ratio} is the number of people above a retirement age,
usually 65, per person of working age, usually 15 to 64.
\end{definition}

In a pay-as-you-go system that balances, contributions equal pensions: $\tau \times \text{wage} \times \text{workers} = \text{pension} \times \text{pensioners}$, so
\begin{equation}\label{eq:b2:overlapping-generations:payg}
\tau = \frac{\text{pensioners}}{\text{workers}} \times \frac{\text{pension}}{\text{wage}},
\end{equation}
the contribution rate equals the dependency ratio times the \omterm{def:b2:overlapping-generations:replacement-rate}{replacement rate}. With 30 pensioners per 100 workers and a \omterm{def:b2:overlapping-generations:replacement-rate}{replacement rate} of one half, contributions take 15 per cent of wages. When the ratio rises to 45 per 100, either the contribution rises to 22.5 per cent, or the \omterm{def:b2:overlapping-generations:replacement-rate}{replacement rate} falls to one third,
or people retire later, which moves people from the numerator to the denominator.

\begin{measure}[The old-age dependency ratio]
The ratio in the opening chart counts people, not workers and pensioners. Not everyone aged 15 to 64 works, some people
work past 65, and many people under 65 receive pensions. A pension system's own ratio of pensioners to contributors can
differ much from the demographic ratio, and moves with employment and the retirement age as well as with births and
deaths. The demographic ratio is still the force behind it: it is set by birth rates decades ago and by how long people
now live, and policy cannot change it quickly.
\end{measure}

\section{Debt across generations}\label{sec:b2:overlapping-generations:debt}

Government debt in this model works like a \omterm{def:b2:overlapping-generations:payg}{pay-as-you-go pension}. A government that borrows from the young and pays the
old, and then rolls the debt over, hands the first old generation a gift and makes the young hold government bonds in
place of capital, so the capital stock falls. If the economy is \omterm{def:b2:overlapping-generations:dynamic-inefficiency}{dynamically inefficient}, with $R < 1 + n$, the debt ratio shrinks of its
own accord as the economy grows, as \cref{ch:b2:government-debt} showed for $i < \gamma$: the debt can be rolled over for ever, and the crowding out of capital
moves the economy toward the golden rule, making every generation better off. In an efficient economy, with $R > 1 + n$, the
debt must eventually be serviced by taxes, and the gift to the first generation is paid for by those that follow.

\section{Exercises}\label{sec:b2:overlapping-generations:exercises}

\begin{exercise}[$\star$]\label{exo:b2:overlapping-generations:1}
Read the opening chart. Which country's ratio rose most since 1960, and which did not rise at all? What drives the
difference?
\end{exercise}

\begin{exercise}[$\star$]\label{exo:b2:overlapping-generations:2}
A balanced pay-as-you-go system has 40 pensioners per 100 workers and a \omterm{def:b2:overlapping-generations:replacement-rate}{replacement rate} of one half. What is the
contribution rate?
\end{exercise}

\begin{exercise}[$\star$]\label{exo:b2:overlapping-generations:3}
Read the efficiency chart. What does a point above the $45^\circ$ line mean, and what fraction of economies lie above it?
\end{exercise}

\begin{exercise}[$\star$]\label{exo:b2:overlapping-generations:4}
What observation in an economy's accounts, sustained over many years, would indicate that it is \omterm{def:b2:overlapping-generations:dynamic-inefficiency}{dynamically
inefficient}?
\end{exercise}

\begin{exercise}[$\star\star$]\label{exo:b2:overlapping-generations:5}
With $\alpha = 1/3$, $\beta = 1$, $n = 0$ and $A = 1$, compute $k^*$ and show that the economy is exactly at the golden rule.
\end{exercise}

\begin{exercise}[$\star\star$]\label{exo:b2:overlapping-generations:6}
Predict, then compute: with $\alpha = 0.3$, is the steady state dynamically efficient when $\beta = 1$? When $\beta = 0.6$? In the inefficient case, who gains from introducing a
small pay-as-you-go system?
\end{exercise}

\begin{exercise}[$\star\star$]\label{exo:b2:overlapping-generations:7}
\textbf{(D)} Using the 2025 dependency ratios of the opening chart and a \omterm{def:b2:overlapping-generations:replacement-rate}{replacement rate} of one half, compute the
contribution rate a balanced pay-as-you-go system would need in Japan, Italy, Korea, the United States and Nigeria. What
does the calculation leave out?
\end{exercise}

\begin{exercise}[$\star\star$]\label{exo:b2:overlapping-generations:8}
Population grows 1 per cent a year and wages per worker 1.5 per cent. What annual return does a pay-as-you-go system pay? Compare
it with a return on capital of 4 per cent.
\end{exercise}

\begin{exercise}[$\star\star$]\label{exo:b2:overlapping-generations:9}
Show that the slope of \eqref{eq:b2:overlapping-generations:k} at the steady state is $\alpha$, and explain what this implies for the speed of convergence.
\end{exercise}

\begin{exercise}[$\star\star\star$]\label{exo:b2:overlapping-generations:10}
Derive \eqref{eq:b2:overlapping-generations:k} from the young person's problem and the firm's first-order conditions.
\end{exercise}

\begin{exercise}[$\star\star\star$]\label{exo:b2:overlapping-generations:11}
Show that in the steady state $R^*/(1 + n) = \alpha(1 + \beta)/(\beta(1 - \alpha))$, and find the condition on $\alpha$ and $\beta$ for \omterm{def:b2:overlapping-generations:dynamic-inefficiency}{dynamic inefficiency}.
\end{exercise}

\begin{exercise}[$\star\star\star$]\label{exo:b2:overlapping-generations:12}
Show that a \omterm{def:b2:overlapping-generations:funded}{funded pension}, with contributions $d$ invested in capital at the return $R$, leaves the capital stock unchanged if
$d \le s_t$.
\end{exercise}

\section{Problem: Brenland's pension reform}\label{sec:b2:overlapping-generations:problem}

\begin{problem}[{Weekend problem --- Brenland's pension reform: the economy without pensions, a pay-as-you-go system, the
generation that gains and the generations that pay}]
\label{pb:b2:overlapping-generations:1}
Brenland is the economy of the text: $\alpha = 0.3$, $\beta = 0.6$, $n = 0.2$, $A = 10$, one period a generation of thirty years; amounts are thousands of dalens per
young person.

\medskip
\noindent\textbf{Part I --- The model.}
\begin{enumerate}
  \item Write the wage as a function of $k$.
  \item Write the gross return on capital.
  \item What share of the wage do the young save?
  \item Write the law of motion of $k$.
  \item Why does the young's saving get divided by $1 + n$?
\end{enumerate}

\medskip
\noindent\textbf{Part II --- No pensions.}
\begin{enumerate}[resume]
  \item Find $k^*$.
  \item Find the wage and the gross return.
  \item What return per year does $R^*$ represent over thirty years?
  \item Find saving, consumption when young and consumption when old.
  \item Find the golden-rule capital stock. Is Brenland dynamically efficient?
  \item Compare the annual return on capital with the annual growth of population.
\end{enumerate}

\medskip
\noindent\textbf{Part III --- A pay-as-you-go system.} Each young person pays $d = 1$.
\begin{enumerate}[resume]
  \item What does each old person receive?
  \item What gross return does a contribution earn?
  \item Using the listing, find the new $k^*$.
  \item Find the new wage and gross return.
  \item Why does capital fall?
  \item Would a funded system with $d = 1$ change $k^*$?
\end{enumerate}

\medskip
\noindent\textbf{Part IV --- Who gains and who pays.}
\begin{enumerate}[resume]
  \item What does the old generation alive when the system starts receive, and what did it pay?
  \item In the new steady state, how much is a young person's contribution worth to her, net, compared with saving it at the return on
    capital?
  \item Why cannot the system be ended without loss to some generation?
  \item Brenland's actual system has 30 pensioners per 100 workers and a \omterm{def:b2:overlapping-generations:replacement-rate}{replacement rate} of one half. What contribution rate balances
    it?
  \item In 2050 there will be 45 pensioners per 100 workers. What contribution rate is needed?
  \item What \omterm{def:b2:overlapping-generations:replacement-rate}{replacement rate} keeps the contribution rate unchanged?
  \item If a later retirement age brought the ratio down to 35 per 100, what contribution rate would balance the system?
  \item State the result: capital per young person with and without the pay-as-you-go system, and the windfall of the first
    old generation, in dalens.
\end{enumerate}
\end{problem}
